% Fragment LaTeX: sieci splotowe (CNN) oraz czasowe sieci splotowe (TCN)
% Wymaga: \usepackage[utf8]{inputenc}, \usepackage[polish]{babel},
%          \usepackage{amsmath}, \usepackage[numbers]{natbib}
% Bibliografia: sieci_splotowe_i_tcn.bib
%
% Kompilacja przykładowa (z katalogu latex/):
%   pdflatex main_sieci_splotowe_i_tcn.tex
%   bibtex main_sieci_splotowe_i_tcn
%   pdflatex main_sieci_splotowe_i_tcn.tex
%   pdflatex main_sieci_splotowe_i_tcn.tex

\section{Sieci splotowe i czasowe sieci splotowe}
\label{sec:cnn-tcn}

Sieci splotowe (\emph{Convolutional Neural Networks}, CNN) oraz czasowe sieci splotowe
(\emph{Temporal Convolutional Networks}, TCN) stanowią rodzinę modeli wykorzystujących
operację splotu do ekstrakcji lokalnych wzorców z danych. CNN historycznie rozwinięto
dla obrazów i sygnałów o strukturze przestrzennej~\citep{lecun1998gradient,lecun2015deep},
natomiast TCN dostosowują splot do modelowania sekwencji i szeregów
czasowych~\citep{lea2017temporal,bai2018empirical}.

%==============================================================================
\subsection{Sieci splotowe (CNN)}
\label{sec:cnn}
%==============================================================================

%Definicja
\subsubsection{Definicja}
Sieć splotowa to architektura uczenia głębokiego, w której podstawową operacją jest
splot dyskretny filtrów (jąder) z lokalnymi fragmentami wejścia~\citep{lecun1998gradient}.
W przeciwieństwie do w pełni połączonych warstw, wagi filtrów są współdzielone
(\emph{weight sharing}) wzdłuż osi przestrzennych lub czasowych, co ogranicza liczbę
parametrów i wprowadza indukcję lokalności oraz ekwiwariancji względem
przesunięć~\citep{goodfellow2016deep}. Typowa CNN składa się z warstw splotowych,
nie\-li\-nio\-wo\-ści (np.\ ReLU), opcjonalnego podpróbkowania (\emph{pooling}) oraz warstw
gęstych na wyjściu~\citep{lecun2015deep}.

%Sposób działania
\subsubsection{Sposób działania}
Dla wejścia $X$ i jądra $K$ o rozmiarze $k$ jednowymiarowy splot w pozycji $t$
można zapisać jako
\begin{equation}
  (X * K)(t) = \sum_{i=0}^{k-1} K(i)\, X(t-i),
  \label{eq:conv1d}
\end{equation}
przy odpowiednim dopełnieniu (\emph{padding}) i kroku (\emph{stride}).
W dwóch wymiarach (obrazy) splot obejmuje obie osie przestrzenne.
Warstwa uczy się wielu filtrów, z których każdy wykrywa inny wzorzec lokalny
(krawędzie, tekstury, motywy czasowe). Głębsze warstwy łączą te cechy w
reprezentacje o większym polu recepcyjnym~\citep{lecun2015deep,goodfellow2016deep}.
Uczenie odbywa się zwykle metodą wstecznej propagacji błędu i stochastycznego
spadku gradientu~\citep{lecun1998gradient}.

%Przykłady
\subsubsection{Przykłady}
Klasyczne zastosowania CNN obejmują rozpoznawanie pisma i dokumentów~\citep{lecun1998gradient}
oraz klasyfikację obrazów w dużej skali~\citep{krizhevsky2012imagenet}.
W domenie szeregów czasowych pełnosplotowe sieci 1D (FCN) stanowią silny punkt
wyjścia do klasyfikacji bez ręcznej inżynierii cech~\citep{wang2017time}.
Do rekonstrukcji i prognozowania przy brakach danych stosowano m.in.\ hybrydy
CNN--LSTM na szeregach zużycia energii~\citep{hussain2022cnnlstm}
oraz architektury oparte na splotach 2D nad wariacjami okresowymi
(TimesNet), obejmujące także zadanie imputacji~\citep{wu2023timesnet}.

%Zalety
\subsubsection{Zalety}
\begin{itemize}
  \item Efektywna ekstrakcja lokalnych wzorców dzięki współdzieleniu wag
        i indukcji lokalności~\citep{lecun2015deep,goodfellow2016deep}.
  \item Znacznie mniej parametrów niż porównywalna sieć w pełni połączona
        przy podobnej mocy ekspresji lokalnej~\citep{lecun1998gradient}.
  \item Dobrze skaluje się na dane o strukturze siatki (obrazy, spektrogramy,
        okna czasowe) i wspiera równoległe obliczenia~\citep{krizhevsky2012imagenet}.
  \item Przenośność filtrów: wczesne warstwy często uczą się uniwersalnych cech,
        co ułatwia transfer uczenia~\citep{yosinski2014transferable}.
\end{itemize}

%Ograniczenia
\subsubsection{Ograniczenia}
\begin{itemize}
  \item Ograniczone pole recepcyjne w płytkich sieciach utrudnia modelowanie
        bardzo długich zależności bez pogłębiania architektury lub użycia
        dilatacji~\citep{yu2016multi}.
  \item Standardowy splot nie jest z natury przyczynowy: bez specjalnego
        dopełnienia może wykorzystywać przyszłe próbki, co jest niedopuszczalne
        w prognozowaniu online~\citep{bai2018empirical}.
  \item Słabsza indukcja dla danych bez lokalnej struktury przestrzennej/czasowej
        (np.\ dowolne zbiory cech tabularnych)~\citep{goodfellow2016deep}.
  \item Interpretacja poszczególnych filtrów bywa trudna, zwłaszcza w głębokich
        stosach warstw~\citep{lecun2015deep}.
\end{itemize}

%==============================================================================
\subsection{Czasowe sieci splotowe (TCN)}
\label{sec:tcn}
%==============================================================================

%Definicja
\subsubsection{Definicja}
Czasowa sieć splotowa (TCN) to architektura sekwencyjna oparta na splotach
przyczynowych i dilatowanych (\emph{causal dilated convolutions}), zaprojektowana
do modelowania zależności czasowych bez rekurencji~\citep{lea2017temporal,bai2018empirical}.
W ujęciu Bai i~wsp.~\cite{bai2018empirical} TCN łączy: (i)~splot przyczynowy
(wyjście w chwili $t$ zależy wyłącznie od wejść $\le t$), (ii)~dilatację zwiększającą
pole recepcyjne wykładniczo wraz z głębokością oraz (iii)~residualne bloki
stabilizujące trening głębokich stosów. Koncepcja dilatowanych splotów przyczynowych
była wcześniej kluczowa m.in.\ w modelu WaveNet~\citep{oord2016wavenet}.

%Sposób działania
\subsubsection{Sposób działania}
Splot dilatowany z czynnikiem $d$ i jądrem rozmiaru $k$ ma postać
\begin{equation}
  (X *_d K)(t) = \sum_{i=0}^{k-1} K(i)\, X(t - d\cdot i),
  \label{eq:dilated}
\end{equation}
gdzie $d$ pomija co $(d-1)$-szą próbkę, rozszerzając zasięg bez zwiększania liczby
parametrów~\citep{yu2016multi,oord2016wavenet}. Przyczynowość uzyskuje się przez
dopełnienie z lewej strony sekwencji, tak aby filtr nie sięgał w
przyszłość~\citep{bai2018empirical}. Typowy blok TCN stosuje dwie warstwy
splotu dilatowanego z normalizacją, nie\-li\-nio\-wo\-ścią i dropoutem oraz połączenie
residualne; dilatacje rosną zwykle według $d = 2^{\ell}$ w kolejnych warstwach
$\ell$, co daje pole recepcyjne rzędu $O(k\cdot 2^{L})$ dla $L$ warstw.
Całą sekwencję można przetwarzać równolegle, w przeciwieństwie do klasycznych
RNN/LSTM~\citep{bai2018empirical}.

%Przykłady
\subsubsection{Przykłady}
TCN i pokrewne architektury stosowano do segmentacji działań w
wideo~\citep{lea2017temporal}, generowania audio~\citep{oord2016wavenet}
oraz ogólnego modelowania sekwencji, gdzie empirycznie dorównywały lub przewyższały
RNN na wielu benchmarkach~\citep{bai2018empirical}.
W prognozowaniu szeregów czasowych z brakami danych wykorzystano m.in.\ grafowe
rozszerzenie splotu czasowego (BiTGraph) z modułem PartialTCN uwzględniającym
wzorce braków~\citep{chen2024bitgraph}.
Pokrewna linia --- TimesNet --- modeluje zmiany czasowe przez sploty 2D na
tensorach okresowych i obejmuje m.in.\ imputację~\citep{wu2023timesnet}.

%Zalety
\subsubsection{Zalety}
\begin{itemize}
  \item Stabilne i zwykle szybsze uczenie niż RNN dzięki pełnej równoległości
        po długości sekwencji~\citep{bai2018empirical}.
  \item Elastyczne, rosnące pole recepcyjne przez dilatacje bez eksplozji
        liczby parametrów~\citep{yu2016multi,oord2016wavenet}.
  \item Ścisła przyczynowość ułatwia uczciwą ewaluację w zadaniach
        prognostycznych~\citep{bai2018empirical}.
  \item Residualne bloki poprawiają przepływ gradientu w głębokich
        sieciach~\citep{he2016deep,bai2018empirical}.
\end{itemize}

%Ograniczenia
\subsubsection{Ograniczenia}
\begin{itemize}
  \item Aby objąć bardzo długie zależności, potrzeba odpowiednio głębokiej sieci
        lub dużych dilatacji, co zwiększa koszt pamięci aktywacji~\citep{bai2018empirical}.
  \item Stałe jądra splotowe gorzej adaptują się do silnie niestacjonarnych
        lub nieregularnie próbkowanych szeregów bez dodatkowych mechanizmów
        (maski, interpolacja, modele grafowe)~\citep{chen2024bitgraph}.
  \item Wybór rozmiaru jądra, schematu dilatacji i głębokości jest zadaniem
        strojenia hiperparametrów o istotnym wpływie na jakość~\citep{bai2018empirical}.
  \item Czyste TCN nie modelują wprost niepewności ani rozkładów braków danych;
        w zadaniach imputacji często wymaga się rozszerzeń lub połączeń
        z innymi rodzinami modeli~\citep{hussain2022cnnlstm,chen2024bitgraph}.
\end{itemize}

\paragraph{Podsumowanie.}
CNN dostarczają ogólnego mechanizmu lokalnej ekstrakcji cech przez splot,
natomiast TCN specjalizują ten mechanizm pod sekwencje: przyczynowość, dilatacje
i residualne bloki pozwalają efektywnie modelować zależności czasowe bez
rekurencji~\citep{lecun2015deep,bai2018empirical}. W rekonstrukcji i prognozowaniu
szeregów czasowych obie linie --- często w wariantach hybrydowych --- stanowią
ważną alternatywę wobec sieci rekurencyjnych i transformerów.
